% Fixture for PDF link layer / Theorem Lens tests. Rebuild with:
%   cd tests/fixtures/pdf/src && pdflatex -interaction=nonstopmode links-paper.tex \
%     && pdflatex -interaction=nonstopmode links-paper.tex && cp links-paper.pdf ..
\documentclass[twocolumn]{article}
\usepackage[a4paper,margin=2cm]{geometry}
\usepackage{hyperref}
\hypersetup{pdftitle={Link Fixture}}
\begin{document}
\section{Introduction}\label{sec:intro}
Sorting was analysed exhaustively by Knuth~\cite{knuth} and later
surveyed by Sedgewick~\cite{sedgewick}. Hash tables follow
Cormen~\cite{clrs}.\footnote{Footnote target text about hashing.}
See Figure~\ref{fig:box}, Equation~\ref{eq:euler} and
Section~\ref{sec:method}. Project page: \url{https://example.org/theorem}.

\begin{equation}\label{eq:euler}
e^{i\pi} + 1 = 0
\end{equation}

\begin{figure}[h]
\centering
\rule{3cm}{2cm}
\caption{A black box figure.}\label{fig:box}
\end{figure}

\section{Method}\label{sec:method}
The method section body text sits here.

\clearpage
\begin{thebibliography}{9}
\bibitem{knuth} Donald E. Knuth. The Art of Computer Programming,
Volume 3: Sorting and Searching. Addison-Wesley, 1998.
\bibitem{sedgewick} Robert Sedgewick. Algorithms in C. Addison-Wesley, 1990.
\bibitem{clrs} Thomas H. Cormen, Charles E. Leiserson, Ronald L. Rivest
and Clifford Stein. Introduction to Algorithms. MIT Press, 2009.
\end{thebibliography}
\end{document}
